\section{To See God Face to Face}

Benedict XII defines in \emph{Benedictus Deus} (1336) that the souls of the just, once purified, enjoy the immediate vision of the divine essence before the resurrection of their bodies and the general judgment. Their blessedness is already real, and the vision continues without interruption. They are with Christ and with the holy angels. Their bodies await resurrection while their souls already possess the promised vision.

The vision is immediate: God himself is its object. The blessed do not merely contemplate a more excellent created representation in place of God. No creature stands between them and him as a substitute object of vision. We now learn through words, images, signs, and effects. These can truly lead us to God, but the promised end exceeds their manner of mediation. God makes himself directly known. The knowledge that comes through faith is fulfilled in sight, and the hope that awaited possession gives way to possession. Charity endures, for the one whom we have begun to love is present in the fullness of the gift.

This vision is not the sight of a bodily shape. God is spirit, and the divine essence is not an object made visible by moving nearer to it in space. Aquinas distinguishes the intellectual vision of God from the operations of bodily sight and imagination (ST I, q.~12, a.~3). Even in our present life, bodily sight is only part of knowing a person. Truth, fidelity, and love engage the mind and heart through the visible world. God will be known more intimately than anything we now know, although the manner of that knowledge exceeds our present experience.

Aquinas calls the God-given perfection that enables this knowledge the \emph{light of glory}. The light strengthens the creature's intellect; it is not another object replacing God (ST I, q.~12, aa.~2, 5). The intellect needs this strengthening because its own power is finite, while the God it will know is infinite. God enables it to know him directly while preserving the difference between his nature and ours.

The blessed see God without comprehending him exhaustively. To comprehend in this strict sense would mean knowing the infinite divine reality with an infinite adequacy. No created mind possesses such a measure (ST I, q.~12, a.~7). Nevertheless, God is truly present to the blessed; they are not receiving a fragment of him, as though his being could be divided into portions. The limitation belongs to the finite knower. God remains inexhaustible even when the blessed know him directly and are wholly happy in that knowledge.

Nor does the vision make every blessed person omniscient. Aquinas explicitly denies that seeing God means knowing everything that God knows or everything that his power could bring into being (a.~8). The fullness of happiness lies in possessing God, not in acquiring divine sovereignty over all possible knowledge. Some of our present curiosity is generous; some is an anxious desire to be beyond surprise or dependence. Heaven fulfills the hunger for truth while healing the impulse to secure ourselves through control. The limit Aquinas describes belongs to the creature's knowledge, while the happiness comes from God himself.

The Catechism also speaks of finding one's true identity in Christ (CCC 1025). The person who has prayed, repented, served, and suffered is the person whom Christ receives. Grace heals what sin has injured in that person and frees him from the vanity and fear that have distorted his conduct. We shall belong to God as ourselves, made whole by the charity for which he created us.

For a person accustomed to prayer in darkness, these truths give patience. The silence of prayer can seem empty because faith has not become sight. Distraction can make even a short act of recollection difficult. Yet the worth of prayer is not measured by how closely its present feeling resembles glory. To remain with God in faith, return after distraction, and entrust oneself to his word are real acts of communion. Perseverance in this prayer depends upon the God who promises to make himself known and who already gives grace to seek him.

Contemplation of the beatific vision also teaches reverence in theological study. We can understand true distinctions about God without possessing the reality as the blessed possess it. Study becomes spiritually fruitful when clearer knowledge nourishes adoration and obedience. For example, knowing that the vision depends upon God's gift gives us a reason to ask for his help while studying and to give thanks for what we understand. We study a gift we cannot yet experience as the blessed do, trusting the God who promises it.

